\section*{Hoofdstuk \ref{ch:g9:periodic-table-first-look} --- Het periodiek systeem: een eerste blik}

\begin{solution}{exo:g9:periodic-table-first-look:1}
Op volgorde van toenemend \omterm{def:g9:inside-the-atom:atomic-number}{atoomnummer} $Z$.
\end{solution}

\begin{solution}{exo:g9:periodic-table-first-look:2}
Koolstof: \omterm{def:g9:periodic-table-first-look:period}{periode} 2, \omterm{def:g9:periodic-table-first-look:period}{groep} 14. Magnesium: \omterm{def:g9:periodic-table-first-look:period}{periode} 3, \omterm{def:g9:periodic-table-first-look:period}{groep} 2. Argon:
\omterm{def:g9:periodic-table-first-look:period}{periode} 3, \omterm{def:g9:periodic-table-first-look:period}{groep} 18.
\end{solution}

\begin{solution}{exo:g9:periodic-table-first-look:3}
\omterm{def:g9:periodic-table-first-look:metal}{Metalen}: ijzer, koper, aluminium. \omterm{def:g9:periodic-table-first-look:metal}{Niet-metalen}: zwavel, zuurstof,
chloor.
\end{solution}

\begin{solution}{exo:g9:periodic-table-first-look:4}
Kalium (\ce{K}) onder natrium; broom (\ce{Br}) onder chloor.
\end{solution}

\begin{solution}{exo:g9:periodic-table-first-look:5}
Calcium (\ce{Ca}, $Z = 20$) en stikstof (\ce{N}, $Z = 7$).
\end{solution}

\begin{solution}{exo:g9:periodic-table-first-look:6}
Het staat in dezelfde groep als natrium, \omterm{def:g9:periodic-table-first-look:period}{groep} 1, en de elementen van
een groep gedragen zich hetzelfde.
\end{solution}

\begin{solution}{exo:g9:periodic-table-first-look:7}
Een \omterm{def:g9:periodic-table-first-look:metal}{metaal}: het glanst, geleidt elektrische stroom en vormt een positief
\omterm{def:g9:ions:ion}{ion}.
\end{solution}

\begin{solution}{exo:g9:periodic-table-first-look:8}
Om elementen met vergelijkbare eigenschappen in dezelfde kolom te
houden, ook als hij daarvoor een vakje leeg moest laten. De gaten waren
voorspellingen: de ontbrekende elementen werden later gevonden, met de
eigenschappen die hij had beschreven.
\end{solution}

\begin{solution}{exo:g9:periodic-table-first-look:9}
Koolstof: $\ce{C} = 12$; zuurstof: $\ce{O} = 16$. De gaten ``? $= 68$'' (naast
Al $= 27.4$) en ``? $= 70$'' (naast Si $= 28$).
\end{solution}

\begin{solution}{exo:g9:periodic-table-first-look:10}
\omterm{def:g9:periodic-table-first-look:period}{Periode} 2: 8 elementen (lithium tot neon). \omterm{def:g9:periodic-table-first-look:period}{Periode} 4: 18 elementen
(kalium tot krypton).
\end{solution}

\begin{solution}{exo:g9:periodic-table-first-look:11}
\omterm{def:g9:periodic-table-first-look:period}{Groep} 18. Overal waar niets mag reageren: het argon in gloeilampen en
lasgas, helium in ballonnen (het brandt niet).
\end{solution}

\begin{solution}{exo:g9:periodic-table-first-look:12}
$(27.0 + 114.8)/2 = 70.9$, dicht bij 69.7.
\end{solution}

\begin{solution}{pb:g9:periodic-table-first-look:1}
\textbf{1.} $Z = 32$; \omterm{def:g9:periodic-table-first-look:period}{periode} 4, \omterm{def:g9:periodic-table-first-look:period}{groep} 14.

\textbf{2.} Erboven: silicium. Eronder: tin. Links: gallium. Rechts:
arseen.

\textbf{3.} Een metalloïde.

\textbf{4.} Het staat in dezelfde groep als silicium, en de elementen van
een groep gedragen zich hetzelfde.

\textbf{5.} $(28.1 + 118.7)/2 = 73.4$.

\textbf{6.} $(69.7 + 74.9)/2 = 72.3$.

\textbf{7.} Het gemiddelde van links en rechts (72.3): langs een \omterm{def:g9:periodic-table-first-look:period}{periode}
neemt het atoomgewicht geleidelijk toe, element na element, terwijl het
naar beneden in een groep met grote stappen springt.

\textbf{8.} $72.6 - 70 = 2.6$.

\textbf{9.} $(28.1 + 118.7 + 69.7 + 74.9)/4 = 291.4/4 = 72.85$, dat is
72.9.

\textbf{10.} Het gemiddelde, 72.9, ligt dicht bij zowel de 72 van
Mendelejev als de 72.3 van Winkler (en de huidige 72.6).

\textbf{11.} Een element dat voorspeld was voordat iemand het had gezien,
werd gevonden met de voorspelde eigenschappen: de tabel was niet zomaar
een lijst, ze kon voorspellingen doen.

\textbf{12.} Ongeveer 72.9.
\end{solution}
